% Chapter 7 -- Eigenvalues and diagonalisation
\chapter{Eigenvalues and diagonalisation}\label{ch:eigen}

Most vectors a linear map touches come out rotated and stretched --- but a
few special directions come out merely \emph{scaled}. Those directions are
the map's eigenvectors, and they are the coordinates the map itself would
choose if it could: in the eigenbasis, a matrix is just a diagonal list of
stretch factors. That one observation unlocks cheap powers of $A$, closed
forms for recurrences, steady states of Markov chains, and --- two chapters
from now --- the whole theory of the SVD. After this chapter you can find
eigenvalues and eigenvectors by hand for small matrices, diagonalise when
diagonalisation exists, say precisely when it fails, and use all of it on a
Fibonacci-style recurrence and a Markov chain.

\section{The eigen equation}\label{sec:eig-eq}

\begin{definition}[Eigenvector, eigenvalue]\index{eigenvector}\index{eigenvalue}
A nonzero vector $\mathbf v$ is an \emph{eigenvector} of $A$ with
\emph{eigenvalue} $\lambda$ when
\[
A\mathbf v=\lambda\mathbf v .
\]
The direction is preserved; only the length (and possibly the sign)
changes. $\mathbf v=\mathbf0$ is excluded on purpose --- it would trivialize
the equation for every $\lambda$.
\end{definition}

Figure~\ref{fig:eigendirections} makes the idea visual: most arrows change
direction under $A$; the eigenvectors are the ones that stay on their own
line.

\begin{figure}[htbp]
\centering
\begin{tikzpicture}[scale=0.8, >=stealth]
  \begin{scope}
    \draw[step=1.0,graybg!50] (-1,-1) grid (3,3);
    \draw[->,gray] (-1.3,0)--(3.3,0); \draw[->,gray] (0,-1.3)--(0,3.3);
    \draw[->,accent,very thick] (0,0)--(1,1) node[above]{\tiny $\mathbf v_1$};
    \draw[->,accent,very thick] (0,0)--(1,-0.5) node[below]{\tiny $\mathbf v_2$};
    \draw[->,gray,dashed] (0,0)--(0.6,0.9) node[left=-1pt]{\tiny $\mathbf w$};
    \node at (1,-1.6) {\small before};
  \end{scope}
  \draw[->,very thick,ink] (4,1)--(5.4,1) node[midway,above]{$A$};
  \begin{scope}[xshift=6.6cm]
    \draw[step=1.0,graybg!50] (-1,-2) grid (6,6);
    \draw[->,gray] (-1.3,0)--(6.3,0); \draw[->,gray] (0,-1.3)--(0,6.3);
    \draw[->,accent,very thick] (0,0)--(5,5) node[above]{\tiny $5\mathbf v_1$};
    \draw[->,accent,very thick] (0,0)--(2,-1) node[below]{\tiny $2\mathbf v_2$};
    \draw[->,gray,dashed] (0,0)--(3.0,2.7);
    \node at (2.5,-1.6) {\small after: same lines, different lengths};
  \end{scope}
\end{tikzpicture}

\caption{Eigenvectors are the map's preserved directions. Notice the two
accent vectors stay on their lines --- merely stretched by $\lambda_1=5$
and $\lambda_2=2$ --- while a generic vector changes direction.
Regenerate: \ttfamily figures/ch07\_eigendirections.tex.}
\label{fig:eigendirections}
\end{figure}

Two immediate facts worth internalizing:

\begin{itemize}
\item Eigenvectors are determined only \emph{up to nonzero scale} --- if
$\mathbf v$ works so does $7\mathbf v$, and $-\mathbf v$. The eigenLINE is
the object; we pick a unit representative when convenient.
\item The pair $(\lambda,\mathbf v)$ still satisfies
$A^2\mathbf v=\lambda^2\mathbf v$, $A^k\mathbf v=\lambda^k\mathbf v$: powers
of $A$ act trivially on eigendirections. This single fact powers the whole
chapter.
\end{itemize}

\section{Characteristic polynomial and finding them by hand}\label{sec:eig-char}

Rewrite $A\mathbf v=\lambda\mathbf v$ as $(A-\lambda I)\mathbf v=\mathbf0$.
A nonzero $\mathbf v$ exists exactly when $A-\lambda I$ is singular ---
chapter 6 says that means $\det(A-\lambda I)=0$.

\begin{definition}[Characteristic polynomial]\index{characteristic polynomial}
$\chi_A(\lambda)=\det(A-\lambda I)$ is a degree-$n$ polynomial in
$\lambda$; its roots are the eigenvalues of $A$.
\end{definition}

\begin{example}[A complete $2\times2$ by hand]\label{ex:eig2}
\textbf{Strategy.} Compute $\chi_A(\lambda)$, factor it (the trace and
determinant give a free checksum: roots sum to $\operatorname{tr}A$ and
multiply to $\det A$), then solve $(A-\lambda I)\mathbf v=\mathbf0$ per
root.

\textbf{Solution.} $A=\begin{bmatrix}4&1\\2&3\end{bmatrix}$:
\[
\det\begin{bmatrix}4-\lambda&1\\2&3-\lambda\end{bmatrix}
=(4-\lambda)(3-\lambda)-2=\lambda^2-7\lambda+10=(\lambda-5)(\lambda-2).
\]
Checksum: $\operatorname{tr}A=7=5+2$, $\det A=10=5\cdot2$. \checkmark
For $\lambda=5$: $A-5I=\begin{bmatrix}-1&1\\2&-2\end{bmatrix}$, null
direction $(1,1)$: indeed $A(1,1)^{T}=(5,5)^{T}$. For $\lambda=2$:
$A-2I=\begin{bmatrix}2&1\\2&1\end{bmatrix}$, direction $(1,-2)$:
$A(1,-2)^{T}=(2,-4)^{T}$. Both verified by direct multiplication.
\end{example}

\begin{lstlisting}[caption={\ttfamily code/ch07\_eig\_by\_hand.py --- the
hand answers checked by residual, plus \ttfamily np.linalg.eig and the two
checksums.}]
import numpy as np
np.set_printoptions(precision=4, suppress=True)

A = np.array([[4., 1.], [2., 3.]])
w_hand, V_hand = 5.0, np.array([1., 1.])
w2_hand, V2_hand = 2.0, np.array([1., -2.])
print("A =")
print(A)
print("hand: lambda=5, v=(1,1):", np.allclose(A @ V_hand, 5 * V_hand))
print("hand: lambda=2, v=(1,-2):", np.allclose(A @ V2_hand, 2 * V2_hand))

w, V = np.linalg.eig(A)
print("np.linalg.eig eigenvalues:", np.sort(w))
print("eig vectors (cols):")
print(V)
print("trace == sum lambda:", np.isclose(np.trace(A), w.sum()))
print("det == prod lambda:", np.isclose(np.linalg.det(A), w.prod()))
\end{lstlisting}

\begin{lstlisting}[caption={Output of the eigen check. Notice the library's
vectors are the same \emph{lines} normalized --- $(1,1)/\sqrt2$ and
$(1,-2)/\sqrt5$ up to sign --- and both checksums pass.}]
A =
[[4. 1.]
 [2. 3.]]
hand: lambda=5, v=(1,1): True
hand: lambda=2, v=(1,-2): True
np.linalg.eig eigenvalues: [2. 5.]
eig vectors (cols):
[[ 0.7071 -0.4472]
 [ 0.7071  0.8944]]
trace == sum lambda: True
det == prod lambda: True
\end{lstlisting}

\begin{exercise}[easy]
Find the eigenvalues of $\begin{bmatrix}3&2\\2&3\end{bmatrix}$ and one
eigenvector for each.
\begin{solution}
$\lambda^2-6\lambda+5=0$: $\lambda=5$ with $\mathbf v=(1,1)$, $\lambda=1$
with $\mathbf v=(1,-1)$.
\end{solution}
\end{exercise}

\begin{exercise}[medium]
Give a $2\times2$ real matrix with \emph{no} real eigenvalues, and compute
its complex pair.
\begin{solution}
Any rotation by $\theta\ne0,\pi$, e.g.
$\begin{bmatrix}0&-1\\1&0\end{bmatrix}$: $\chi=\lambda^2+1$, roots
$\pm i$. A real matrix \emph{can} have complex eigenvalues --- they arrive
in conjugate pairs (Figure~\ref{fig:complex-rot}).
\end{solution}
\end{exercise}

\begin{figure}[htbp]
\centering
\begin{tikzpicture}[scale=0.8, >=stealth]
  \draw[step=1.0,graybg!50] (-2.5,-2.5) grid (2.5,2.5);
  \draw[->,gray] (-2.8,0)--(2.8,0); \draw[->,gray] (0,-2.8)--(0,2.8);
  \draw[->,accent,very thick] (0,0)--(1.8,0) node[below right=-2pt]{\tiny $\mathbf v$};
  \draw[->,amber,very thick] (0,0)--(1.27,1.27) node[right]{\tiny $R_{45^\circ}\mathbf v$};
  \draw[->,amber,thick] (0,0)--(0,1.8) node[above]{\tiny $R_{45^\circ}^2\mathbf v$};
  \draw[->,gray] (1.5,0.3) arc (15:60:0.8);
  \node at (2.1,1.0) {\tiny $45^\circ$};
  \node at (0,-3.3) {\small rotation $\leftrightarrow$ complex pair};
\end{tikzpicture}

\caption{Complex eigenvalues mean rotation. Notice a $45^\circ$ rotation
has no real preserved line at all --- its eigenvalues are
$e^{\pm i\pi/4}$. Regenerate: \ttfamily figures/ch07\_complex\_rotation.tex.}
\label{fig:complex-rot}
\end{figure}

\section{Eigenspaces and multiplicities}\label{sec:eig-mult}

\begin{definition}[Eigenspace]\index{eigenspace}
$E_\lambda=\ker(A-\lambda I)$ is the \emph{eigenspace} of $\lambda$ --- all
eigenvectors for that $\lambda$ plus $\mathbf0$. Its dimension is the
\emph{geometric multiplicity}\index{geometric multiplicity}; the number of
times $\lambda$ repeats as a root of $\chi_A$ is the \emph{algebraic
multiplicity}\index{algebraic multiplicity}. Always
$\dim E_\lambda\le(\text{algebraic multiplicity})$.
\end{definition}

Example: $\begin{bmatrix}1&1\\0&1\end{bmatrix}$ has $\lambda=1$ with
algebraic multiplicity 2, but $E_1=\operatorname{span}\{(1,0)\}$ is only
1-dimensional --- a defective matrix (Section~\ref{sec:eig-defective}).

\begin{theorem}\label{th:eig-indep}
Eigenvectors belonging to \emph{distinct} eigenvalues are linearly
independent.
\end{theorem}

\begin{proof}
Suppose $\mathbf v_1,\dots,\mathbf v_k$ are eigenvectors for distinct
$\lambda_1,\dots,\lambda_k$ and some combination vanishes. Apply $A$
repeatedly: $c_1\lambda_1^j\mathbf v_1+\cdots+c_k\lambda_k^j\mathbf v_k=0$
for all $j$. Pick two consecutive $j$ and eliminate $\mathbf v_k$: the
remaining combination has coefficients $c_i(\lambda_i-\lambda_k)$, nonzero
only where $\lambda_i\ne\lambda_k$ --- induction kills them one at a time.
\end{proof}

Consequence: $n$ distinct eigenvalues $\Rightarrow$ a full eigenbasis
$\Rightarrow$ diagonalisable. The converse is false in both directions
(repeated eigenvalues can still diagonalise; defective ones cannot).

\section{Diagonalisation and powers of \texorpdfstring{$A$}{A}}\label{sec:eig-diag}

Stack the eigenvectors as columns of $V$ and the eigenvalues into
$\Lambda=\operatorname{diag}(\lambda_1,\dots,\lambda_n)$:
$AV=V\Lambda$ column-by-column. When $V$ is invertible --- exactly when $A$
has $n$ independent eigenvectors ---

\begin{theorem}[Diagonalisation]\index{diagonalisation}
$A=V\Lambda V^{-1}$ iff $A$ has a basis of eigenvectors, and then
\[
A^k=V\Lambda^{k}V^{-1},\qquad
\Lambda^k=\operatorname{diag}(\lambda_1^k,\dots,\lambda_n^k).
\]
\end{theorem}

This is the change of basis of chapter 4 paying its second dividend: the
same map, in its own basis, is a list of independent scalings.

\begin{example}[$A^{10}$ the cheap way]\label{ex:a10}
\textbf{Strategy.} Don't multiply ten matrices --- decompose the input in
the eigenbasis, scale each piece by $\lambda^{10}$, reassemble.

\textbf{Solution.} From Example~\ref{ex:eig2},
$A=V\Lambda V^{-1}$ with $V=\begin{bmatrix}1&1\\1&-2\end{bmatrix}$,
$\Lambda=\operatorname{diag}(5,2)$. Take $\mathbf x=(1,0)$: solving
$V\mathbf c=\mathbf x$ gives $\mathbf c=(2/3,\,1/3)$, i.e.
$\mathbf x=\tfrac23\mathbf v_1+\tfrac13\mathbf v_2$. Then
\[
A^{10}\mathbf x=\tfrac23\,5^{10}\mathbf v_1+\tfrac13\,2^{10}\mathbf v_2
=\begin{bmatrix}(2\cdot5^{10}+2^{10})/3\\(2\cdot5^{10}-2\cdot2^{10})/3\end{bmatrix}
=\begin{bmatrix}6510758\\6509734\end{bmatrix}.
\]
Two exponentiations instead of nine matrix multiplies --- and the dominant
term $5^{10}$ already explains the answer's scale.
\end{example}

\begin{lstlisting}[caption={\ttfamily code/ch07\_power\_iteration.py ---
eight steps of power iteration: multiply, normalize, read the Rayleigh
quotient.}]
import numpy as np
np.set_printoptions(precision=6, suppress=True)

A = np.array([[2., 1., 0.], [1., 2., 1.], [0., 1., 2.]])
x = np.array([1., 0.5, -0.25])
print("A =")
print(A)
for k in range(8):
    x = A @ x
    x = x / np.linalg.norm(x)
    rq = x @ A @ x
    print(f"k={k+1}: lambda~{rq:.6f}  x={x}")
w, V = np.linalg.eig(A)
print("true dominant:", np.abs(w).max(), " eigenvector:", V[:, np.argmax(np.abs(w))])
\end{lstlisting}

\begin{lstlisting}[caption={Output of power iteration. Notice the iterate
rotates onto the dominant eigendirection and the Rayleigh quotient locks
onto $\lambda_1=2+\sqrt2\approx3.4142$ --- the sign flip versus
\texttt{np.linalg.eig} is the up-to-scale freedom, not an error.}]
A =
[[2. 1. 0.]
 [1. 2. 1.]
 [0. 1. 2.]]
k=1: lambda~2.939597  x=[0.819232 0.573462 0.      ]
k=2: lambda~3.205318  x=[0.733761 0.652232 0.190234]
k=3: lambda~3.334830  x=[0.653366 0.686872 0.318307]
k=4: lambda~3.385931  x=[0.594985 0.699983 0.39499 ]
k=5: lambda~3.404379  x=[0.557226 0.704641 0.439295]
k=6: lambda~3.410823  x=[0.534021 0.706258 0.46478 ]
k=7: lambda~3.413048  x=[0.52009  0.706815 0.479498]
k=8: lambda~3.413814  x=[0.511822 0.707007 0.488037]
true dominant: 3.4142135623730914  eigenvector: [-0.5      -0.707107 -0.5     ]
\end{lstlisting}

\begin{figure}[htbp]
\centering
\begin{tikzpicture}[scale=0.8, >=stealth]
  \draw[step=1.0,graybg!50] (-1,-1) grid (5,5);
  \draw[->,gray] (-1.3,0)--(5.3,0); \draw[->,gray] (0,-1.3)--(0,5.3);
  \draw[gray,dashed] (0,0)--(5,3.54);
  \node[below right=-2pt] at (4.6,3.0) {\tiny dominant line};
  \draw[->,accent!30,thick] (0,0)--(1.0,0.1);
  \draw[->,accent!55,thick] (0,0)--(2.0,0.9);
  \draw[->,accent!80,thick] (0,0)--(3.2,2.0);
  \draw[->,accent,very thick] (0,0)--(4.4,3.1);
  \node at (2.5,-0.8) {\small iterates rotate onto the dominant eigendirection};
\end{tikzpicture}

\caption{Power iteration geometrically: each application leans harder on
the largest eigenvalue, so iterates rotate onto the dominant
eigendirection. This is the algorithm inside PageRank and every
large-scale eigensolver. Regenerate:
\ttfamily figures/ch07\_power\_convergence.tex.}
\label{fig:power-conv}
\end{figure}

\section{Defective matrices}\label{sec:eig-defective}

When algebraic multiplicity outgrows geometric multiplicity, the
eigenbasis runs out of members and $V$ cannot be inverted.

\begin{example}[A shear that cannot diagonalise]\label{ex:defective}
\textbf{Strategy.} Compute both multiplicities for the shear
$S=\begin{bmatrix}1&1\\0&1\end{bmatrix}$ and look at the geometry directly.

\textbf{Solution.} $\chi_S=(1-\lambda)^2$: $\lambda=1$ with algebraic
multiplicity 2. But $S-I=\begin{bmatrix}0&1\\0&0\end{bmatrix}$ has null
space $\operatorname{span}\{(1,0)\}$ --- geometric multiplicity 1. Only the
$x$-axis survives as a direction (Figure~\ref{fig:defective}); the second
slot in the would-be eigenbasis is missing, so $S\ne V\Lambda V^{-1}$ for
any $V$. The correct replacement (Jordan form) adds a \emph{generalized}
eigenvector: $S$ maps $\mathbf e_2$ to $\mathbf e_2+\mathbf e_1$, and
powers show it: $S^k=\begin{bmatrix}1&k\\0&1\end{bmatrix}$ grows
\emph{linearly}, not exponentially --- the signature of a defect.
\end{example}

\begin{figure}[htbp]
\centering
\begin{tikzpicture}[scale=0.8, >=stealth]
  \begin{scope}
    \draw[step=1.0,graybg!50] (-1,-1) grid (3,3);
    \draw[->,gray] (-1.3,0)--(3.3,0); \draw[->,gray] (0,-1.3)--(0,3.3);
    \draw[->,accent,very thick] (0,0)--(1.5,0) node[below]{\tiny $\mathbf v$};
    \draw[->,good,very thick] (0,0)--(0.8,1.5) node[above]{\tiny $\mathbf w$};
    \node at (1,-1.6) {\small before: two directions};
  \end{scope}
  \draw[->,very thick,ink] (4,1)--(5.4,1) node[midway,above]{\tiny shear};
  \begin{scope}[xshift=6.6cm]
    \draw[step=1.0,graybg!50] (-1,-1) grid (4,3);
    \draw[->,gray] (-1.3,0)--(4.3,0); \draw[->,gray] (0,-1.3)--(0,3.3);
    \draw[->,accent,very thick] (0,0)--(1.5,0) node[below]{\tiny $\mathbf v$};
    \draw[->,good,very thick] (0,0)--(2.3,1.5) node[above]{\tiny $\mathbf w+\mathbf v$};
    \draw[dashed,gray] (0.8,1.5)--(2.3,1.5);
    \node at (1.5,-1.6) {\small after: only $\mathbf v$ survives as a line};
  \end{scope}
\end{tikzpicture}

\caption{The shear's eigendirections collapse. Notice only the
$x$-direction is preserved; every other direction slides along it and the
second eigenvector simply does not exist. Regenerate:
\ttfamily figures/ch07\_defective.tex.}
\label{fig:defective}
\end{figure}

\begin{exercise}[medium]
Decide whether $A=\begin{bmatrix}2&0\\0&2\end{bmatrix}$ and
$B=\begin{bmatrix}2&1\\0&2\end{bmatrix}$ diagonalise, though both have
$\chi=(2-\lambda)^2$.
\begin{solution}
$A=2I$: every vector is an eigenvector, basis exists, already diagonal.
$B$: $\ker(B-2I)=\operatorname{span}\{(1,0)\}$, dimension 1 $<2$ ---
defective, does not diagonalise.
\end{solution}
\end{exercise}

\begin{exercise}[hard]
Show that if $\lambda$ is an eigenvalue of $A$ with eigenvector
$\mathbf v$, then for every polynomial $p$, $p(\lambda)$ is an eigenvalue
of $p(A)$ with the same eigenvector. Conclude the eigenvalues of $A^2$ are
$\lambda_i^2$ --- and explain why $A^2$ can have \emph{fewer} distinct
eigenvalues than $A$.
\begin{solution}
$p(A)\mathbf v=\sum_j c_jA^j\mathbf v=\sum_j c_j\lambda^j\mathbf v
=p(\lambda)\mathbf v$. For $p(x)=x^2$: eigenvalues $\lambda_i^2$; if
$\pm\lambda$ both occur they merge into one $\lambda^2$ (e.g.\
eigenvalues $2$ and $-2$ of $A$ both give $4$ in $A^2$).
\end{solution}
\end{exercise}

\section{Applications: recurrences and Markov chains}\label{sec:eig-apps}

\subsection*{Fibonacci, closed form}

The recurrence $F_{k+1}=F_k+F_{k-1}$ becomes a matrix iteration:
\[
\begin{bmatrix}F_{k+1}\\F_{k}\end{bmatrix}
=\begin{bmatrix}1&1\\1&0\end{bmatrix}
\begin{bmatrix}F_{k}\\F_{k-1}\end{bmatrix},
\qquad\text{i.e.}\quad \mathbf x_{k}=A\mathbf x_{k-1}=A^{k}\mathbf x_0 .
\]

\begin{example}[Fibonacci via diagonalisation]\label{ex:fib}
\textbf{Strategy.} Eigenvalues of $A$: $\chi_A=\lambda^2-\lambda-1$,
roots $\varphi=(1+\sqrt5)/2$ and $\psi=(1-\sqrt5)/2$. Distinct $\Rightarrow$
diagonalisable; the full closed form is derived in
Section~\ref{sec:handderiv7}.

\textbf{Solution.} With $F_0=0,F_1=1$ the eigencoordinate computation
yields Binet's formula,
\[
F_n=\frac{\varphi^{n}-\psi^{n}}{\sqrt5},
\]
and since $|\psi|<1$ its term dies exponentially: $F_n$ is the nearest
integer to $\varphi^{n}/\sqrt5$.
\end{example}

\begin{lstlisting}[caption={\ttfamily code/ch07\_fibonacci.py --- Binet's
formula checked against the recursion, timed, and pushed past its float
limit by an integer matrix power.}]
import numpy as np
import time
np.set_printoptions(precision=4, suppress=True)

A = np.array([[1., 1.], [1., 0.]])
w, V = np.linalg.eig(A)
print("eigenvalues of [[1,1],[1,0]]:", np.sort(w))
phi, psi = (1 + np.sqrt(5)) / 2, (1 - np.sqrt(5)) / 2

def fib_closed(n):
    return int(round((phi**n - psi**n) / np.sqrt(5)))

def fib_rec(n):
    a, b = 0, 1
    for _ in range(n):
        a, b = b, a + b
    return a

GOOD = [n for n in range(1, 100) if fib_closed(n) == fib_rec(n)]
print("closed form exact through n =", GOOD[-1],
      "; float64 Binet breaks at n =", GOOD[-1] + 1)
for n in [10, 20, 30, 60, 70]:
    assert fib_closed(n) == fib_rec(n), n
print("closed form == recursion for n in {10,20,30,60,70}: OK")

t0 = time.perf_counter()
for n in range(1, 70):
    fib_closed(n)
t1 = time.perf_counter()
for n in range(1, 70):
    fib_rec(n)
t2 = time.perf_counter()
print(f"closed-form all n<70: {1e3*(t1-t0):.3f} ms")
print(f"recursion   all n<70: {1e3*(t2-t1):.3f} ms")
Aint = np.array([[1, 1], [1, 0]], dtype=object)
Ai = np.linalg.matrix_power(Aint, 100)
print("F_101 via integer matrix power:", Ai[0, 0])
n = 10
An = np.linalg.matrix_power(A, n)
print(f"A^{n} =")
print(An.astype(int), f"(entries F11={int(An[0,0])}, F10={int(An[1,0])})")
\end{lstlisting}

\begin{lstlisting}[caption={Output of the Fibonacci check. Notice float64
Binet holds only to $n=70$ --- past that the $\varphi^{n}$ term outgrows
precision; integer matrix powers stay exact forever. (Timing is a wash at
this $n$; both are cheap --- the point is the \emph{form} exists at all.)}]
eigenvalues of [[1,1],[1,0]]: [-0.618  1.618]
closed form exact through n = 70 ; float64 Binet breaks at n = 71
closed form == recursion for n in {10,20,30,60,70}: OK
closed-form all n<70: 0.076 ms
recursion   all n<70: 0.047 ms
F_101 via integer matrix power: 573147844013817084101
A^10 =
[[89 55]
 [55 34]] (entries F11=89, F10=55)
\end{lstlisting}

\subsection*{Markov chains and the eigenvalue 1}

A Markov chain stores a probability distribution $\mathbf p$ and updates
$\mathbf p\leftarrow P\mathbf p$ where each column of $P$ sums to 1.
Columns summing to 1 makes $\lambda=1$ an eigenvalue, and every steady
state is its eigenvector:

\begin{example}[Weather chain steady state]\label{ex:markov}
\textbf{Strategy.} For the two-state chain
$P=\begin{bmatrix}0.9&0.3\\0.1&0.7\end{bmatrix}$ (columns: from-state),
find the $\lambda=1$ eigenvector and check that powers converge to it.

\textbf{Solution.} $(P-I)\boldsymbol\pi=\mathbf0$:
$\begin{bmatrix}-0.1&0.3\\0.1&-0.3\end{bmatrix}\boldsymbol\pi=\mathbf0$
gives $\pi_1=3\pi_2$; normalizing $\pi_1+\pi_2=1$ yields
$\boldsymbol\pi=(0.75,0.25)$. The second eigenvalue is $0.6<1$, so every
other eigendirection decays: $P^{k}\mathbf x\to\boldsymbol\pi$ regardless
of start --- visible in the listing (Figure~\ref{fig:markov-graph} shows
the chain).
\end{example}

\begin{figure}[htbp]
\centering
\begin{tikzpicture}[scale=1.0, >=stealth]
  \node[circle,draw=accent,very thick,minimum size=1.1cm] (R) at (0,0) {rain};
  \node[circle,draw=accent,very thick,minimum size=1.1cm] (S) at (4,0) {sun};
  \draw[->,thick] (R) edge[loop above] node{$0.9$} ();
  \draw[->,thick] (S) edge[loop above] node{$0.7$} ();
  \draw[->,thick] (R) edge[bend left=25] node[above]{$0.1$} (S);
  \draw[->,thick] (S) edge[bend left=25] node[below]{$0.3$} (R);
  \node at (2,-1.3) {\small $\pi=(0.75,\,0.25)$: eigenvector of $\lambda=1$};
\end{tikzpicture}

\caption{The weather chain's transition graph. Notice the steady state is
an \emph{eigenvector}: it is what the dynamics settle onto, not a separate
concept. Regenerate: \ttfamily figures/ch07\_markov\_graph.tex.}
\label{fig:markov-graph}
\end{figure}

\begin{lstlisting}[caption={\ttfamily code/ch07\_markov\_steady.py ---
steady state from \ttfamily np.linalg.eig, from matrix powers, and from a
Monte-Carlo run of 200{,}000 steps.}]
import numpy as np
np.set_printoptions(precision=6, suppress=True)
np.random.seed(1)

P = np.array([[0.9, 0.3],
              [0.1, 0.7]])
print("P =")
print(P)
print("columns sum to 1:", np.allclose(P.sum(axis=0), 1))

w, V = np.linalg.eig(P)
print("eigenvalues:", w)
pi = V[:, np.argmax(w)].real
pi = pi / pi.sum()
print("steady state (eigenvector of lambda=1):", pi)

x = np.array([1., 0.])
for k in [10, 20, 50, 200]:
    xk = np.linalg.matrix_power(P, k) @ x
    print(f"P^{k} x0 = {xk}")
rng = np.random.default_rng(0)
state, counts = 0, np.zeros(2)
for _ in range(200000):
    state = rng.choice(2, p=P[:, state])
    counts[state] += 1
print("Monte Carlo fractions:", counts / counts.sum())
\end{lstlisting}

\begin{lstlisting}[caption={Output of the Markov check. Notice three
routes to the same answer --- eigenvector $(0.75,0.25)$, powers converging
by $k=20$, Monte-Carlo fractions within sampling error. The subdominant
eigenvalue $0.6$ sets the convergence speed.}]
P =
[[0.9 0.3]
 [0.1 0.7]]
columns sum to 1: True
eigenvalues: [1.  0.6]
steady state (eigenvector of lambda=1): [0.75 0.25]
P^10 x0 = [0.751512 0.248488]
P^20 x0 = [0.750009 0.249991]
P^50 x0 = [0.75 0.25]
P^200 x0 = [0.75 0.25]
Monte Carlo fractions: [0.753955 0.246045]
\end{lstlisting}

The same $\lambda=1$-eigenvector machinery is PageRank in disguise: the
``random surfer'' matrix is column-stochastic (with a teleport term to
kill dangling nodes), and a page's rank is its steady-state probability ---
computed by power iteration, exactly the loop above, at web scale
(Figure~\ref{fig:dominant-geo} is the picture).

\begin{figure}[htbp]
\centering
\begin{tikzpicture}[scale=0.75, >=stealth]
  \draw[step=1.0,graybg!50] (-1,-1) grid (6,5);
  \draw[->,gray] (-1.3,0)--(6.3,0); \draw[->,gray] (0,-1.3)--(0,5.3);
  \draw[->,accent,very thick] (0,0)--(1.4,0.7) node[below right=-2pt]{\tiny $\mathbf v_1\ \lambda_1{=}3$};
  \draw[->,good,very thick] (0,0)--(0.4,1.5) node[above left=-2pt]{\tiny $\mathbf v_2\ \lambda_2{=}0.5$};
  \draw[->,amber,very thick] (0,0)--(4.2,2.1) node[below right=-2pt]{\tiny $A\mathbf v_1=3\mathbf v_1$};
  \draw[->,good!60,thick,dashed] (0,0)--(0.05,0.19) node[left=-2pt]{\tiny $A^3\mathbf v_2$};
  \node at (3,-1.0) {\small dominant $\lambda_1$ grows, subdominant $\lambda_2$ dies};
\end{tikzpicture}

\caption{Why the dominant eigenvalue rules: the $\lambda_1$ component is
tripled per application while the $\lambda_2$ component shrinks by half.
After a few steps only the dominant direction is visible.
Regenerate: \ttfamily figures/ch07\_dominant\_geometry.tex.}
\label{fig:dominant-geo}
\end{figure}

\section{Hand derivation}\label{sec:handderiv7}

Derive Binet's formula from the eigenvalues. With
$A=\begin{bmatrix}1&1\\1&0\end{bmatrix}$,
$\chi_A=\lambda^2-\lambda-1$, roots
$\varphi,\psi=(1\pm\sqrt5)/2$. Eigenvectors: $(\lambda,1)^{T}$ works since
row 2 of $(A-\lambda I)$ reads $1-\lambda\cdot0$... directly:
$A(\lambda,1)^{T}=(\lambda+1,\lambda)^{T}=(\lambda^2,\lambda)^{T}
=\lambda(\lambda,1)^{T}$ using $\lambda^2=\lambda+1$. So
$V=\begin{bmatrix}\varphi&\psi\\1&1\end{bmatrix}$. Write the start:
$(1,0)^{T}=c_1\mathbf v_\varphi+c_2\mathbf v_\psi$ --- solving
$c_1\varphi+c_2\psi=1$, $c_1+c_2=0$ gives $c_1=-c_2=1/(\varphi-\psi)=1/\sqrt5$.
Then
\[
\mathbf x_n=A^{n}\mathbf x_0=\frac{\varphi^{n}}{\sqrt5}\mathbf v_\varphi
-\frac{\psi^{n}}{\sqrt5}\mathbf v_\psi,
\qquad
F_n=\frac{\varphi^{n}-\psi^{n}}{\sqrt5}.
\]
Check at $n=1$: $(\varphi-\psi)/\sqrt5=1$ \checkmark; at $n=2$:
$(\varphi^2-\psi^2)/\sqrt5=(\varphi+\psi)(\varphi-\psi)/\sqrt5=1\cdot1=1$
\checkmark\ (using $\varphi+\psi=1$, $\varphi\psi=-1$ from
$\lambda^2=\lambda+1$).

\section{Common mistakes}

\begin{tripwire}[{Expecting real matrices to have real eigenvalues}]
A rotation has none. Real matrices get complex eigenvalues in conjugate
pairs --- and they are not a pathology, they \emph{are} the rotation
content of the map (Figure~\ref{fig:complex-rot}).
\end{tripwire}

\begin{tripwire}[{Sign errors in $\det(A-\lambda I)$}]
For $2\times2$, $\det(A-\lambda I)=(a-\lambda)(d-\lambda)-bc$; the
cross-term is a \emph{minus} $bc$ and the $\lambda$ sits on the diagonal,
not the off-diagonal. Guard yourself with the free checksum
$\sum\lambda_i=\operatorname{tr}A$, $\prod\lambda_i=\det A$.
\end{tripwire}

\begin{tripwire}[{Assuming every matrix diagonalises}]
Distinct eigenvalues guarantee it; repeated eigenvalues are ambiguous.
$\begin{bmatrix}1&1\\0&1\end{bmatrix}$ is the canonical counterexample ---
check $\dim E_\lambda$, not just the algebraic count.
\end{tripwire}

\begin{tripwire}[{Treating eigenvectors as unique}]
They are directions, defined only up to nonzero scale and sign.
\texttt{np.linalg.eig} returns unit vectors whose signs may disagree with
your hand work --- compare directions ($|\cos\theta|=1$ or
\texttt{allclose} up to sign), never entries literally.
\end{tripwire}

\begin{tripwire}[{Confusing eigenvalues of $A$ and $A^2$}]
$A^2$ has eigenvalues $\lambda_i^2$ with the \emph{same} eigenvectors ---
not $\lambda_i$ again, and not unrelated values. Opposite eigenvalues
($\pm\lambda$) collide in $A^2$, so $A^2$ can have fewer distinct values.
\end{tripwire}

\section{Summary}

\begin{center}
\begin{tabularx}{\textwidth}{l>{\raggedright\arraybackslash}X}
\toprule
\textbf{Object} & \textbf{What you need to know}\\
\midrule
Eigen equation & $A\mathbf v=\lambda\mathbf v$, $\mathbf v\ne\mathbf0$;
direction preserved \\
Characteristic poly & $\chi_A(\lambda)=\det(A-\lambda I)$; roots $=$
eigenvalues; $\operatorname{tr}=\sum\lambda_i$, $\det=\prod\lambda_i$ \\
Eigenspace & $E_\lambda=\ker(A-\lambda I)$; geometric $\le$ algebraic
multiplicity \\
Independence & eigenvectors of distinct $\lambda_i$ are linearly
independent \\
Diagonalisation & $A=V\Lambda V^{-1}$ iff $n$ independent eigenvectors;
$A^k=V\Lambda^kV^{-1}$ \\
Defective & geometric $<$ algebraic $\Rightarrow$ no eigenbasis (shear);
powers grow polynomially \\
Power iteration & multiply + normalize $\to$ dominant eigenpair;
rate set by $|\lambda_2/\lambda_1|$ \\
Markov steady state & column-stochastic $P$ has $\lambda_1=1$; steady
state is its eigenvector \\
\bottomrule
\end{tabularx}
\end{center}

\section*{Additional exercises}
\addcontentsline{toc}{section}{Additional exercises}

\begin{exercise}[easy]
Find eigenvalues and eigenvectors of
$\begin{bmatrix}2&0\\0&-3\end{bmatrix}$ and
$\begin{bmatrix}1&2\\0&-1\end{bmatrix}$. What does a zero row on the
diagonal buy you?
\end{exercise}

\begin{exercise}[easy]
$\mathbf v$ is an eigenvector of $A$ with $\lambda=4$. What are the
eigenvalues of $A^{3}$, $A^{-1}$, $A+2I$ on the same vector?
\end{exercise}

\begin{exercise}[easy]
Compute $\chi_A$ for $A=\begin{bmatrix}1&3\\0&5\end{bmatrix}$ and confirm
$\operatorname{tr}=\sum\lambda_i$, $\det=\prod\lambda_i$.
\end{exercise}

\begin{exercise}[easy]
Give the algebraic and geometric multiplicities of every eigenvalue of
$\operatorname{diag}(2,2,3)$.
\end{exercise}

\begin{exercise}[medium]
Diagonalise $A=\begin{bmatrix}1&2\\2&1\end{bmatrix}$ and compute
$A^8(1,0)^{T}$ exactly via the eigen decomposition.
\end{exercise}

\begin{exercise}[medium]
For the Markov matrix $P=\begin{bmatrix}0.8&0.4\\0.2&0.6\end{bmatrix}$,
find the steady state and the subdominant eigenvalue; estimate how many
steps until a start at $(1,0)$ is within $0.01$ of steady state.
\end{exercise}

\begin{exercise}[medium]
Show that similar matrices have the same eigenvalues: if $B=P^{-1}AP$ and
$A\mathbf v=\lambda\mathbf v$, then $B(P^{-1}\mathbf v)=\lambda(P^{-1}\mathbf v)$.
Does the converse --- same eigenvalues means similar --- hold?
\end{exercise}

\begin{exercise}[medium]
In code: run power iteration on a symmetric $4\times4$ of your choice,
then on one whose two largest eigenvalues are close ($|\lambda_1/\lambda_2|\approx1.05$).
Count iterations to $10^{-6}$ Rayleigh-quotient error in each case and
relate to $|\lambda_2/\lambda_1|^k$.
\end{exercise}

\begin{exercise}[medium]
The matrix $J=\begin{bmatrix}1&1&0\\0&1&1\\0&0&1\end{bmatrix}$. Compute
$J^4$ exactly --- the polynomial growth of a defective matrix shows up as
binomial coefficients.
\end{exercise}

\begin{exercise}[hard]
Prove: $A$ diagonalises iff the sum of geometric multiplicities equals
$n$ iff the minimal polynomial has no repeated roots. (State the minimal
polynomial; you may use $\chi_A$'s factorization and the eigenspace
decomposition.)
\end{exercise}

\begin{exercise}[hard]
Analyse $R_\theta=\begin{bmatrix}\cos\theta&-\sin\theta\\\sin\theta&\cos\theta\end{bmatrix}$:
find its (complex) eigenvalues and eigenvectors, show
$R_\theta=VDV^{-1}$ over $\mathbb C$, and explain why diagonalisation over
$\mathbb R$ fails for $\theta\ne0,\pi$.
\end{exercise}

\begin{exercise}[hard]
Implement shifted power iteration: to find the eigenvalue \emph{nearest}
$\mu$, run power iteration on $(A-\mu I)^{-1}$. Test on a $3\times3$ of
your choice against \texttt{np.linalg.eig}, and explain why the dominant
eigenvalue of $(A-\mu I)^{-1}$ is $1/(\lambda_i-\mu)$.
\end{exercise}

\begin{exercise}[hard]
A page graph with dangling node: build the $3\times3$ matrix
$P=\begin{bmatrix}0&0.5&1/3\\0.5&0&1/3\\0.5&0.5&1/3\end{bmatrix}$ plus a
teleport term $0.85P+0.15\cdot\frac13J$ ($J$ all-ones). Find the PageRank
vector by power iteration in code and confirm $\lambda_1=1$.
\end{exercise}

\section*{Notes}
\addcontentsline{toc}{section}{Notes}

\chapternotes{%
The eigen-equation-first development and the applications section follow
Strang \cite{ila} ch.~10 and 18.06SC Unit II \cite{ocw1806sc}; linear
dynamics as the motivating application follows VMLS ch.~9 \cite{vmls} and
MML \S4.4 \cite{mml}. Axler ch.~5 \cite{axler} develops eigen theory
\emph{without} determinants (invariant subspaces first) --- the opposite
order from this book's characteristic-polynomial route, worth reading as a
contrast once the computation here is comfortable.}
